% !TEX root = ../meu-trabalho.tex
% -----------------------------------------------------------------------------
% Referencial teórico
% -----------------------------------------------------------------------------

\chapter{Referencial teórico}
\label{chap_referencial_teorico}

\section{Tomada de decisão no mercado financeiro}
\label{sec_tomada_de_decisao_no_mercado_financeiro}

\citeonline{neugebauer2026portfolio} aponta que, especialmente em épocas de incerteza, é fundamental estabelecer uma estratégia para combater fenômenos financeiros adversos, a exemplo da inflação, protegendo assim a renda contra a degradação. Para guiar essa tomada de decisão, os fatores mais comumente analisados pelos indivíduos são o retorno esperado, o risco e a liquidez do investimento. Além do caráter protetivo, \citeonline{adil2022financial} ressaltam que os investidores estão frequentemente dispostos a realizar novos aportes impulsionados pelo desejo de obter melhores retornos financeiros, muito embora seja necessária cautela durante essas operações. Com o acesso facilitado a diversas vias e instrumentos, a utilização desses fundos de forma racional, buscando a maximização do retorno aliada à minimização do risco, constitui o objetivo central na alocação de recursos \cite{paiva2019decision,adil2022financial}.

O cenário macroeconômico também exerce um papel decisivo na motivação daqueles que buscam alocar seu capital. Conforme explicam \citeonline{umar2021bitcoin}, o aumento das incertezas em relação às políticas econômicas ou conflitos políticos torna os investidores pessimistas no que diz respeito ao ambiente internacional. Como consequência dessa instabilidade, as pessoas são compelidas a procurar proteção, direcionando seus recursos para ativos que ofereçam segurança, a fim de compensar as eventuais perdas geradas pela incerteza do mercado \cite{umar2021bitcoin}.

Apesar de essas motivações buscarem o lucro ou a proteção, as razões pelas quais as pessoas investem esbarram em um debate central sobre a racionalidade e a previsibilidade dos ganhos. A teoria financeira tradicional, fundamentada na Hipótese de Mercado Eficiente (HME), pressupõe que os investidores agem de forma puramente lógica e que os preços refletem integralmente todas as informações disponíveis, o que tornaria impossível a obtenção de retornos anormais consistentes \cite{paiva2019decision}. A literatura mais recente, contudo, tem contestado essa visão, argumentando que os mercados não operam de forma totalmente aleatória ou eficiente, visto que as decisões de investimento são frequentemente influenciadas por fatores psicológicos, vieses emocionais e limitações cognitivas que podem levar a escolhas irracionais \cite{almansour2023behavioral}. Desse modo, ao constatar que as séries financeiras possuem memória e padrões de comportamento identificáveis, os indivíduos também investem porque se torna possível delinear modelos de previsibilidade para buscar lucros extraordinários e diversificar portfólios, contrariando as premissas da eficiência de mercado \cite{paiva2019decision,charfeddine2020investigating}.

De acordo com as teorias financeiras convencionais, o investimento ocorre sob a premissa de que os indivíduos são agentes racionais, capazes de processar e avaliar toda a informação existente no mercado antes de decidir \cite{adil2022financial}. \citeonline{paiva2019decision} complementam que, nessa visão tradicional, o processo decisório é centralizado na análise do binômio risco e retorno esperado, o que leva o investidor a selecionar combinações de ativos que garantam a menor variância possível para um lucro desejado ou o maior retorno para um nível de risco aceitável. Essa lógica fundamenta a busca pela eficiência estatística na montagem de portfólios, tratando a alocação de recursos como um problema de otimização matemática para que se alcance o melhor equilíbrio entre as opções disponíveis \cite{paiva2019decision}.

Entretanto, as finanças comportamentais revelam que os investidores frequentemente se distanciam da racionalidade pura, investindo de forma influenciada por estados psicológicos e emocionais que provocam divergências em relação ao comportamento previsto pela teoria clássica \cite{adil2022financial}. Conforme argumentam \citeonline{almansour2023behavioral}, o comportamento do mercado é afetado por uma gama de fatores subjetivos, incluindo vieses cognitivos, influências sociais e percepções de risco que mudam dinamicamente conforme o contexto. Essa instabilidade na percepção individual e a tendência de seguir o comportamento de grupos podem induzir a escolhas de investimento de baixa qualidade, resultando em desempenhos de carteira inferiores aos modelos puramente lógicos \cite{almansour2023behavioral,adil2022financial}.

Como resposta à natureza caótica e complexa dos mercados, tem se consolidado uma abordagem de investimento baseada na busca por padrões históricos que possam indicar movimentos futuros de preços \cite{paiva2019decision}. Para \citeonline{paiva2019decision}, embora o mercado seja extremamente sensível a fatores externos, o reconhecimento de padrões de comportamento nas séries financeiras permite delinear modelos de previsibilidade mais robustos. Assim, os métodos atuais de investimento evoluem para sistemas que utilizam ferramentas de aprendizado de máquina para classificar ativos com maior potencial de ganho, integrando esses sinais tecnológicos a modelos de diversificação clássicos para definir de forma mais precisa a alocação de capital e mitigar os efeitos da irracionalidade humana \cite{paiva2019decision}.

No contexto da diversificação de carteiras e da busca por retornos ajustados ao risco, a consolidação das criptomoedas representa a emersão de uma nova classe de ativos financeiros. Esse crescimento acelerado tem atraído a atenção de diferentes perfis de investidores interessados em gerar lucros adicionais e explorar novas oportunidades de composição de portfólio \cite{charfeddine2020investigating}. Do ponto de vista estrutural, ativos como o Bitcoin operam como moedas digitais descentralizadas, funcionando de forma independente das instituições financeiras tradicionais e de autoridades governamentais centralizadas \cite{umar2021bitcoin}.

Essa independência tecnológica é fundamentada na integração de princípios criptográficos e protocolos de rede de ponta a ponta, os quais asseguram um mecanismo de pagamento tolerante a falhas e altamente seguro para transações eletrônicas \cite{ahmed2020risk}. A inovação técnica por trás dessa dinâmica, viabilizada pela tecnologia blockchain, consiste na substituição da figura do intermediário central pela validação descentralizada, transferindo a garantia de confiança da operação para uma rede de computadores que atuam em consenso \cite{charfeddine2020investigating}. Como resultado, essa estrutura impulsionou as criptomoedas a um status de instrumento de investimento global, absorvendo o escrutínio de gestores, reguladores e formuladores de políticas \cite{ahmed2020risk}.

Apesar de compartilharem essa base tecnológica, o ecossistema de criptoativos apresenta uma heterogeneidade significativa que afeta o comportamento dos investimentos. Enquanto criptomoedas convencionais possuem alta volatilidade e operam sob dinâmicas próprias de precificação, existem moedas projetadas para manter estabilidade, conhecidas como stablecoins \cite{goodell2021diversifying}. Ativos como o Tether, por exemplo, possuem seu valor indexado a moedas fiduciárias, como o dólar norte-americano, apresentando características e utilidades inerentemente distintas de ativos como o Bitcoin \cite{goodell2021diversifying}. Essa diversidade intrínseca à classe exige que a inclusão de criptomoedas em estratégias de investimento seja analisada por meio de modelos robustos, dado que cada tipo de moeda digital pode exercer um papel diferente na alocação de capital e na exposição ao risco.

A inclusão de criptomoedas em carteiras de investimento é amplamente justificada na literatura pela sua capacidade de proporcionar diversificação. Diversos estudos apontam que ativos digitais, como o Bitcoin, apresentam uma correlação historicamente fraca ou até mesmo são desconectados em relação aos mercados financeiros tradicionais, incluindo índices de ações, metais preciosos e commodities de energia \cite{ahmed2020risk,charfeddine2020investigating}. Essa falta de sincronia ocorre porque as criptomoedas não compartilham necessariamente dos mesmos determinantes de preço que os ativos convencionais, estando menos expostas aos indicadores macroeconômicos clássicos \cite{charfeddine2020investigating}. Desse modo, a alocação de proporções relativamente pequenas de capital nesses ativos pode gerar benefícios significativos para a redução do risco geral do portfólio, configurando uma estratégia eficiente de diversificação \cite{charfeddine2020investigating}.

No que diz respeito ao momento ideal para o investimento, a literatura investiga frequentemente o papel das criptomoedas como um "porto seguro" (safe haven) em períodos de elevada instabilidade. \citeonline{umar2021bitcoin} observam que, diante de picos de incerteza política e econômica, a demanda por Bitcoin tende a aumentar, uma vez que os investidores buscam ativos alternativos para compensar as perdas sofridas nas bolsas tradicionais. No entanto, os autores ressaltam que essa característica de proteção não é linear, operando em ciclos. Há períodos de retração em que os ativos digitais sofrem quedas simultâneas aos instrumentos financeiros convencionais, o que exige do investidor a habilidade de identificar o ciclo de mercado correto antes de adotar a criptomoeda como um instrumento de hedge \cite{umar2021bitcoin}.

A eficácia desse investimento como proteção torna-se ainda mais específica durante crises econômicas extremas, fenômeno que costuma elevar a correlação entre as diferentes classes de ativos. \citeonline{goodell2021diversifying} demonstram que, em choques de grande magnitude como o provocado pela pandemia de COVID-19, criptomoedas convencionais e de alta volatilidade (como Bitcoin, Ethereum e Litecoin) passam a se mover na mesma direção dos índices de ações, falhando em atuar como diversificadores eficientes. Em contrapartida, as stablecoins, a exemplo do Tether, apresentam um comportamento de correlação negativa com as ações durante essas turbulências \cite{goodell2021diversifying}. Portanto, segundo a literatura, o investimento em criptoativos pareados a moedas fiduciárias é recomendado especialmente em cenários de forte retração econômica e restrição de liquidez, atuando como verdadeiros portos seguros e reservas de valor \cite{goodell2021diversifying}.

\section{Cripto Ativos}
\label{sec_cripto_ativos}

\subsection{Conceito, origem e funcionamento dos criptoativos}
\label{subsec_conceito_origem_e_funcionamento_dos_criptoativos}

As criptomoedas são ativos digitais concebidos para operar como meios de troca descentralizados, utilizando funções criptográficas para assegurar a propriedade, conduzir transações financeiras e controlar a emissão de novas unidades \cite{doran2014forensic,ballis2020herding,chevallier2021forecast}. O Bitcoin consolidou-se como a primeira criptomoeda descentralizada, baseada fundamentalmente na tecnologia blockchain \cite{fakhfekh2020volatility}, tendo sido idealizado em 2008 sob o pseudônimo Satoshi Nakamoto, que designa um indivíduo ou grupo anônimo responsável por publicar o documento fundacional do ativo em uma lista de e-mails sobre criptografia \cite{charfeddine2020investigating}. A blockchain atua como a tecnologia estrutural do sistema, operando como um livro-razão digital público e sem necessidade de permissão, no qual as transações entre os usuários são invariavelmente registradas \cite{sebastiao2021forecasting}.

O funcionamento operacional das criptomoedas ocorre por meio de uma arquitetura de pagamentos ponto a ponto (peer-to-peer), permitindo que os fluxos financeiros ocorram diretamente entre as partes envolvidas \cite{wang2020stablecoins}. A segurança desse ecossistema é garantida de forma estrita pela criptografia, eliminando a dependência de fatores humanos ou relações de confiança \cite{narayanan2016bitcoin}. A inovação intrínseca ao Bitcoin e à blockchain reside na transferência do mecanismo de criação de confiança: de uma terceira parte centralizada para um corpo de consenso coletivo formado por computadores independentes na rede \cite{charfeddine2020investigating}. Consequentemente, a característica definidora primordial das criptomoedas é a exclusão absoluta de instituições financeiras atuando como intermediárias nas operações \cite{harwick2016cryptocurrency,wang2020stablecoins}.

As criptomoedas têm se consolidado como um novo instrumento financeiro e uma classe de ativos com características singulares, embora ainda persista ceticismo e falta de compreensão sobre a sua verdadeira natureza no mercado \cite{charfeddine2020investigating,corbet2020contagion}. O mercado de Bitcoin, em particular, estabeleceu-se de forma sólida como um mercado global de ativos \cite{park2023intelligent}.

Apesar de o Bitcoin carecer de valor intrínseco que lhe permita gerar dividendos ou ser amplamente utilizado como meio de pagamento convencional, a moeda obteve amplo reconhecimento como um ativo de investimento. Sua atratividade financeira fundamenta-se na possibilidade de comercialização a preços superiores, uma dinâmica impulsionada pela escassez restritiva de sua oferta digital \cite{park2023intelligent}. Esse processo de legitimação levou fundos de cobertura e gestores de capital a incorporarem criptomoedas em seus portfólios, estimulando também a comunidade acadêmica a direcionar pesquisas para a negociação desses ativos, com ênfase na aplicação de algoritmos de aprendizado de máquina \cite{sebastiao2021forecasting}.

Por fim, constata-se uma heterogeneidade estrutural dentro dessa nova classe de ativos. Existem criptoativos, a exemplo do Tether, que diferem inerentemente do Bitcoin por possuírem seu valor indexado a moedas fiduciárias, como o dólar norte-americano, o que lhes confere utilidades e características financeiras fundamentalmente distintas no mercado \cite{goodell2021diversifying}.

\subsection{Características do mercado de criptoativos}
\label{subsec_caracteristicas_do_mercado_de_criptoativos}

O mercado de criptomoedas, atualmente um dos maiores mercados não regulamentados do mundo \cite{sebastiao2021forecasting}, é estruturalmente caracterizado por sua extrema volatilidade. Os preços do Bitcoin e de outros ativos digitais apresentam oscilações substanciais não apenas em horizontes de longo prazo, mas também em bases diárias \cite{baur2021crypto}. Quando comparadas a classes de ativos convencionais, como ações, títulos de dívida, ouro e moedas fiduciárias, as cotações do Bitcoin chegam a ser dezenas de vezes mais voláteis e exibem uma sensibilidade aguda a eventos de mercado e a anúncios regulatórios \cite{fakhfekh2020volatility,park2023intelligent}.

A natureza de alta volatilidade desse mercado impõe dificuldades práticas para investidores que buscam a obtenção de retornos estáveis ou a manutenção de valor no tempo \cite{wang2020stablecoins}. O ambiente de negociação é marcado por uma dinâmica de alto risco e altas recompensas, evidenciando padrões recorrentes de retornos atípicos e quedas frequentes de preços \cite{park2023intelligent}. O risco inerente aos criptoativos é exacerbado pela falta de recursos legais disponíveis para a proteção dos investidores em casos de fraude ou falhas operacionais \cite{charfeddine2020investigating}. Ademais, a literatura financeira destaca a ocorrência do comportamento de manada (herding behavior) como um fenômeno investigado e característico das dinâmicas de atuação dos participantes nesse mercado \cite{ballis2020herding}.

A inclusão de criptomoedas em portfólios de investimento possui amplo respaldo na literatura acadêmica devido ao seu potencial intrínseco de proporcionar diversificação \cite{charfeddine2020investigating}. Essa capacidade fundamenta-se na constatação de que a correlação cruzada entre os ativos digitais e os ativos financeiros convencionais é fraca, embora apresente variações ao longo do tempo \cite{charfeddine2020investigating}. O Bitcoin, de forma específica, demonstra estar segmentado em relação aos mercados financeiros tradicionais e de commodities, exibindo maior reatividade ao fluxo de informações gerado por outras criptomoedas, o que, consequentemente, afeta a precisão de suas previsões \cite{chevallier2021forecast}.

Internamente ao mercado digital, os retornos das diferentes criptomoedas apresentam alta correlação entre si, além de exibirem elevada exposição a riscos de cauda \cite{fakhfekh2020volatility}. Não obstante, a estruturação de carteiras diversificadas contendo múltiplas criptomoedas possibilita a mitigação desses riscos, gerando desempenhos superiores --- mensurados pelo índice de Sharpe e retornos equivalentes de certeza --- quando comparados à detenção de um ativo digital isolado \cite{brauneis2019cryptocurrency}.

No que tange à proteção em cenários adversos, a utilidade desses ativos possui limitações. Em condições normais de mercado, as stablecoins operam fundamentalmente como diversificadores eficazes, assumindo a função de porto seguro (safe haven) apenas sob circunstâncias muito específicas \cite{wang2020stablecoins}. Em períodos de severa ruptura econômica e financeira, as evidências indicam que os criptoativos perdem sua capacidade de atuar como instrumentos de proteção (hedge) ou portos seguros, podendo funcionar como amplificadores do contágio financeiro sistêmico \cite{corbet2020contagion}.

A literatura especializada destaca que os retornos do Bitcoin não se adequam à Hipótese de Mercado Eficiente (HME), indicando que as variações de preços não são eventos independentes, mas sim passíveis de previsão \cite{fakhfekh2020volatility}. Essa ineficiência manifesta-se de forma heterogênea e decorre, em grande parte, da ausência de regulamentações estruturais quando comparada aos ativos financeiros convencionais \cite{park2023intelligent}. Tais falhas geram uma volatilidade turbulenta, evidenciando que o mercado não opera de modo puramente aleatório e que técnicas preditivas isoladas podem ser insuficientes para lidar com essa dinâmica \cite{park2023intelligent}.

A persistência de ineficiências invalida a premissa de eficiência plena, viabilizando a identificação de padrões para a extração de lucros. Diante dessas oportunidades, observou-se nos últimos anos um aumento expressivo no interesse pelo uso de técnicas de aprendizado de máquina (machine learning) para prever o comportamento das criptomoedas \cite{sebastiao2021forecasting}. Essa tecnologia consolidou-se como um instrumento eficaz para o desenvolvimento de estratégias de negociação, pois possui a capacidade de inferir relações nos dados que são, frequentemente, imperceptíveis à observação humana direta \cite{fang2022cryptocurrency}.

O uso de modelos de aprendizado de máquina proporciona níveis elevados de precisão, melhorando a previsibilidade dos retornos e superando o desempenho de abordagens concorrentes tradicionais, como os modelos de médias móveis exponenciais e médias móveis integradas autorregressivas \cite{sebastiao2021forecasting}. Uma vantagem crucial dessa tecnologia em relação aos modelos econométricos clássicos é que o aprendizado de máquina não exige que os dados sejam estacionários para garantir a obtenção de soluções válidas. Dado que a estacionariedade é uma premissa restritiva e essencial da econometria, mas irrelevante para a inteligência artificial, as técnicas de machine learning despontam como uma via metodológica altamente promissora para o setor \cite{chevallier2021forecast}.

Contudo, a aplicação desses algoritmos esbarra em desafios relacionados à sua explicabilidade. Mesmo apresentando alta acurácia na otimização de portfólios, as ferramentas de aprendizado de máquina operam frequentemente como "caixas pretas", fornecendo previsões precisas sem que as razões lógicas subjacentes à decisão sejam claras para o investidor \cite{babaei2022explainable}. As criptomoedas continuam sendo um ativo complexo para a modelagem matemática, mantendo o debate sobre a precisão e os limites dos métodos de previsão de seus preços ainda em aberto na literatura \cite{chevallier2021forecast}.

\subsection{Limitações e desafios do investimento em criptoativos}
\label{subsec_limitacoes_e_desafios_do_investimento_em_criptoativos}

A alta volatilidade intrínseca às criptomoedas impõe dificuldades substanciais para a obtenção de retornos estáveis e a preservação de valor no mercado \cite{wang2020stablecoins}. Em decorrência dessas oscilações extremas e da falta de regulamentações, os investidores enfrentam obstáculos diretos para estabelecer estratégias preditivas eficazes, visto que ineficiências heterogêneas causam turbulências que técnicas de previsão isoladas não conseguem superar \cite{park2023intelligent}. As decisões financeiras nesse ambiente são frequentemente embasadas em informações limitadas, motivadas por perspectivas de lucros de curto prazo e sujeitas a resultados altamente incertos \cite{fang2022cryptocurrency}. Adicionalmente, \citeonline{corbet2020contagion} indicam que, em períodos de grave ruptura econômica e financeira, esses ativos falham em atuar como instrumentos de proteção (hedge) ou portos seguros, funcionando, na verdade, como amplificadores de contágio sistêmico.

A escassez de regulamentação, quando comparada à de outros ativos financeiros, é um dos fatores centrais para a ocorrência do comportamento de manada e de ineficiências no mercado \cite{park2023intelligent}. Por operarem sem a supervisão de governos ou bancos centrais, esses ativos geram preocupações quanto aos potenciais impactos negativos na estabilidade do sistema monetário e financeiro global \cite{charfeddine2020investigating}. Consequentemente, o mercado apresenta riscos severos, não oferecendo recursos ou garantias legais aos investidores em situações de fraude ou falhas operacionais \cite{charfeddine2020investigating}.

Sob a perspectiva tecnológica, a infraestrutura das criptomoedas enfrenta severas limitações de escalabilidade. \citeonline{fang2022cryptocurrency} destacam que o volume e a velocidade das transações nessas redes ainda não conseguem competir com a negociação de moedas tradicionais, o que já ocasionou acúmulos e atrasos operacionais de múltiplos dias em plataformas de negociação.

No âmbito comportamental, a literatura aponta que os investidores do setor principal de criptomoedas frequentemente agem de forma irracional, mimetizando as decisões de terceiros sem qualquer referência às suas próprias convicções \cite{ballis2020herding}. Esse comportamento de manada faz com que os ativos se movam em conjunto, gerando dinâmicas e tendências que não refletem, necessariamente, os seus verdadeiros fundamentos econômicos subjacentes \cite{ballis2020herding}.

\section{Machine Learning}
\label{sec_machine_learning}

\subsection{Fundamentos do Machine Learning}
\label{subsec_fundamentos_do_machine_learning}

O Aprendizado de Máquina (Machine Learning - ML) compreende uma subárea da inteligência artificial que se destaca por sua capacidade de prever eventos futuros com base na análise de dados históricos (HAMAYEL; OWDA, 2021). O princípio fundamental dessa tecnologia reside na aptidão dos algoritmos para extrair conhecimento e aprender representações internas diretamente a partir de dados brutos (séries temporais, por exemplo), eliminando a necessidade de conhecimento especializado do domínio para a filtragem manual de características de entrada (LIVIERIS et al., 2020).

A evolução dessas técnicas ocorreu em grande parte como resposta às limitações dos métodos estatísticos e econométricos tradicionais, que conseguiam lidar apenas com problemas lineares e frequentemente enfrentavam dificuldades com a multicolinearidade de múltiplas variáveis. Com o avanço tecnológico, modelos de ML como Redes Neurais Artificiais (ANNs), Máquinas de Vetores de Suporte (SVM) e Florestas Aleatórias (Random Forest) passaram a ser adotados justamente por conseguirem analisar e processar relações não lineares altamente complexas entre diversos fatores \cite{park2023intelligent}. Esse amadurecimento permitiu que os algoritmos entregassem previsões que não apenas se igualam à realidade, mas que aumentam substancialmente a precisão dos resultados em comparação com modelos anteriores (HAMAYEL; OWDA, 2021).

Atualmente, a relevância do Machine Learning é evidenciada pela sua aplicação bem-sucedida em inúmeras áreas de pesquisa de sistemas complexos. Conforme destacam Derbentsev et al. (2021), algoritmos de ML têm demonstrado eficácia significativa em campos variados, como biomedicina, neurorobótica, análise de imagem e voz, e reconhecimento de padrões, o que motivou sua recente e forte aplicação na análise de séries temporais financeiras. Nesse contexto contemporâneo, estratégias modernas como o aprendizado profundo (deep learning) e o aprendizado combinado (ensemble learning) alcançam níveis de precisão de ponta, pois são capazes de criar características valiosas enquanto filtram eficientemente o ruído dos dados de entrada (LIVIERIS et al., 2020).

No aprendizado de máquina supervisionado, os problemas são frequentemente categorizados em tarefas de regressão e classificação, dependendo do objetivo específico da previsão a ser realizada \cite{park2023intelligent}. Em aplicações financeiras, a regressão é comumente utilizada para estimar o valor exato de um ativo no futuro, enquanto a classificação é aplicada para prever a direção do movimento dos preços, como identificar se a cotação vai subir ou descer no período seguinte (LIVIERIS et al., 2020; DERBENTSEV et al., 2021). Para garantir a robustez na construção e validação dessas ferramentas em séries temporais, os dados costumam ser divididos em subconjuntos distintos: um conjunto de treinamento, utilizado para ajustar e estimar os parâmetros; um conjunto de validação, para selecionar o modelo mais adequado; e um conjunto de teste, que avalia o desempenho final e a capacidade do modelo de generalizar informações em dados não vistos anteriormente (SEBASTIÃO; GODINHO, 2021; DERBENTSEV et al., 2021).

Diante da complexidade natural das tarefas de regressão e classificação, ocorre com frequência que modelos individuais de aprendizado de máquina não consigam fornecer, de maneira isolada, a precisão e a qualidade preditiva desejadas (DERBENTSEV et al., 2021). Como uma solução sofisticada para esse desafio, o aprendizado combinado (ensemble learning) propõe a integração de múltiplos modelos para reduzir a alta variância dos preditores individuais e minimizar o erro de generalização (LIVIERIS et al., 2020). O princípio fundamental dessa estratégia consiste em agregar as previsões de diversos algoritmos considerados "fracos" ou de baixo viés, fazendo com que os erros específicos de cada modelo sejam compensados mutuamente durante o processo de previsão, o que invariavelmente resulta em uma performance global superior (DERBENTSEV et al., 2021; PARK; YANG, 2023).

Dentro dessa arquitetura de múltiplos modelos, destacam-se duas abordagens principais: o bagging e o boosting \cite{park2023intelligent}. Enquanto no bagging os modelos básicos aprendem de forma independente e paralela, compensando os resultados por meio de uma média ou votação simples, o boosting opera por meio de uma construção sequencial e adaptativa, onde cada modelo atual é desenvolvido com base nos resultados do modelo anterior (DERBENTSEV et al., 2021; PARK; YANG, 2023). Consequentemente, o método de boosting se utiliza de uma votação ponderada em vez de simples, focando seu esforço computacional em corrigir as falhas sucessivas dos classificadores prévios, o que lhe confere um grande potencial para extrair padrões não lineares complexos dos dados (DERBENTSEV et al., 2021).

\subsection{AdaBoost: funcionamento e características}
\label{subsec_adaboost_funcionamento_e_caracteristicas}

O algoritmo AdaBoost, cujo nome deriva de Adaptive Boosting, fundamenta-se na premissa de reduzir o viés do modelo ao combinar diferentes técnicas de aprendizado de máquina em um único sistema de previsão robusto (BUSARI; LIM, 2021). O conceito central dessa abordagem reside na utilização de modelos considerados "fracos", que realizam o aprendizado de maneira sequencial e adaptativa \cite{park2023intelligent}. Na prática, esses classificadores fracos costumam ser árvores de decisão rasas, as quais, apesar de possuírem baixa precisão individualmente, são computacionalmente eficientes e fáceis de projetar (DERBENTSEV et al., 2021).

No funcionamento interno do algoritmo, as observações do conjunto de treinamento são ponderadas de forma que cada novo classificador seja treinado levando em conta o sucesso ou fracasso dos modelos anteriores (BUSARI; LIM, 2021). Após cada etapa de treinamento, os pesos dos dados são redistribuídos: as amostras que foram classificadas incorretamente recebem um peso maior, garantindo que os próximos aprendizes concentrem seus esforços nos casos mais difíceis (BUSARI; LIM, 2021). Intuitivamente, quanto maior o erro cometido por um modelo prévio em um dado específico, maior será a importância atribuída a esse dado na rodada seguinte para tentar corrigir a falha de previsão \cite{park2023intelligent}.

A decisão final do AdaBoost é construída por meio da soma ponderada das saídas de todos os submodelos criados durante as iterações \cite{park2023intelligent}. Esse processo de integração leva em conta a confiança ou credibilidade de cada modelo: quanto maior for a acurácia de um submodelo específico, maior será o peso da sua resposta no resultado \cite{park2023intelligent}. Essa arquitetura de soma de classificadores gera um modelo global fortalecido, que aumenta consideravelmente a força preditiva do sistema e entrega resultados com precisão superior (BUSARI; LIM, 2021).

Dentre as principais vantagens do AdaBoost, destacam-se a sua alta escalabilidade devido à mecânica simples, a exigência de poucos parâmetros de ajuste e a sua eficácia comprovada na minimização do sobre ajuste (overfitting) \cite{park2023intelligent}. Embora o algoritmo consiga reduzir tanto a variância quanto o viés das previsões, ele possui limitações importantes (DERBENTSEV et al., 2021). Conforme apontam \citeonline{park2023intelligent}, o uso de uma função de perda exponencial pode induzir a erros quando a base de dados possui muito ruído, pois o algoritmo tende a superestimar a importância de valores atípicos (outliers), prejudicando a performance final.

\subsection{Aplicações de Machine Learning em finanças e criptoativos}
\label{subsec_aplicacoes_de_machine_learning_em_financas_e_criptoativos}

O mercado de criptoativos, caracterizado por sua natureza dinâmica e alta complexidade, tornou-se um campo de interesse para a aplicação de algoritmos de aprendizado de máquina (HAMAYEL; OWDA, 2021). De acordo com Derbentsev et al. (2021), esses algoritmos apresentam um desempenho superior aos modelos tradicionais de séries temporais na previsão de taxas de câmbio e preços de fechamento diários das moedas digitais mais capitalizadas. Nesse contexto, Livieris et al. (2020) ressaltam que o desenvolvimento de modelos preditivos inteligentes é considerado uma necessidade fundamental para a tomada de decisão estratégica e para a otimização de portfólios em um ambiente financeiro globalizado.

Um dos maiores desafios enfrentados pelos investidores nesse setor é a volatilidade extrema, que chega a ser significativamente superior à observada em mercados convencionais de ações ou ouro \cite{park2023intelligent}. Livieris et al. (2020) explicam que o mercado de criptomoedas é marcado por fortes flutuações de preços ao longo do tempo, o que exige ferramentas capazes de lidar com dados não estacionários para garantir soluções válidas. Além disso, a presença de eventos anômalos e extremos reforça a importância de técnicas robustas de aprendizado profundo para identificar padrões ocultos e mitigar riscos em períodos de incerteza (SEBASTIÃO; GODINHO, 2021; DERBENTSEV et al., 2021).

No âmbito das operações de trading, a previsão precisa dos movimentos de preços permite que investidores financeiros reduzam seus riscos ao estabelecer políticas de investimento mais adequadas (LIVIERIS et al., 2020). Modelos baseados em aprendizado de máquina auxiliam na identificação de momentos ideais para a compra e venda de ativos, integrando informações históricas de preços com indicadores técnicos, como médias móveis (HAMAYEL; OWDA, 2021; DERBENTSEV et al., 2021). Conforme demonstram \citeonline{sebastiao2021forecasting}, a utilização de conjuntos de modelos (ensembles) pode gerar estratégias de negociação lucrativas mesmo sob condições adversas de mercado, superando estratégias passivas tradicionais.

Fatores comportamentais e externos também desempenham um papel crucial na formação de preços dos ativos digitais. \citeonline{park2023intelligent} destacam a ocorrência frequente de comportamentos de manada e ineficiências de mercado, que podem causar turbulências prolongadas devido a quebras estruturais. Diante disso, Hamayel e Owda (2021) apontam que a análise de sentimento extraída de mídias sociais e plataformas de microblogs constitui um fator relevante que pode influenciar diretamente o volume de negociação e a volatilidade do mercado de criptomoedas.

Por fim, a alocação de criptoativos e o uso de consultores automatizados (robot advisors) têm se expandido para tornar os mercados mais transparentes e reduzir custos operacionais \cite{babaei2022explainable}. Entretanto, embora os modelos de aprendizado de máquina apresentem alta acurácia na otimização de carteiras, eles costumam operar como "caixas-pretas", o que dificulta a compreensão das razões lógicas por trás das decisões de alocação \cite{babaei2022explainable}. Para superar essa barreira, \citeonline{park2023intelligent} propõem sistemas de negociação inteligentes que integram conhecimentos econômicos sobre turbulência de mercado para maximizar recompensas e emitir alertas sobre investimentos de alto risco.

\section{Teoria Moderna de Portfólio}
\label{sec_teoria_moderna_de_portfolio}

\subsection{Fundamentos da teoria moderna de portfólio}
\label{subsec_fundamentos_da_teoria_moderna_de_portfolio}

A gestão de portfólio caracteriza-se como um processo analítico voltado à seleção e alocação de um grupo de ativos financeiros, no qual a distribuição dos investimentos é ajustada continuamente com o objetivo de otimizar o retorno esperado e a tolerância ao risco \cite{chaweewanchon2022markowitz}. O alicerce dessa prática reside no modelo de média-variância (MV) proposto originalmente em 1952, que é considerado a base da teoria de portfólio e permanece como uma referência amplamente reconhecida na gestão de investimentos \cite{chaweewanchon2022markowitz}.

O trabalho seminal de Markowitz (1952) é apontado como o marco inicial da teoria financeira moderna, introduzindo uma solução matemática para gerenciar o dilema entre a busca pelo maior retorno possível e a redução do risco associado ao investimento \cite{paiva2019decision}. De acordo com \citeonline{paiva2019decision}, a formulação desse modelo fundamenta-se na estruturação de carteiras de investimento que buscam maximizar as expectativas de ganho ao mesmo tempo em que mitigam a exposição ao risco.

Sob o ponto de vista comportamental e de mercado, \citeonline{jeleskovic2024cryptocurrency} destacam que a Teoria Moderna de Portfólio parte da premissa de que os investidores são racionais e avessos ao risco, buscando composições de carteira que otimizem a relação entre retorno esperado e risco. Nesse contexto, o modelo média-variância de Markowitz avalia os ativos não de forma isolada, mas conforme sua contribuição para o risco e o retorno da carteira, considerando especialmente variância, covariância e diversificação.

\subsection{Modelo média-variância de Markowitz}
\label{subsec_modelo_media_variancia_de_markowitz}

O modelo de média-variância (MV) fundamenta-se na utilização da média e da variância, extraídas de preços históricos, para mensurar o retorno esperado e o risco de uma carteira \cite{chaweewanchon2022markowitz}. \citeonline{bakry2021bitcoin} apontam que essa abordagem estabelece uma estrutura de otimização baseada em objetivos conflitantes, buscando maximizar o retorno enquanto se minimiza a variância, que atua como a métrica representativa do risco.

Dentro dessas premissas conceituais, o risco é definido pela flutuação dos retornos em relação à sua média aritmética \cite{jeleskovic2024cryptocurrency}. Consequentemente, o processo de seleção de ativos por parte do investidor deve priorizar opções que entreguem a menor variância para um nível de retorno específico ou, de forma inversa, o maior retorno esperado para um limite de risco previamente estabelecido \cite{paiva2019decision}.

\subsection{Diversificação e fronteira eficiente}
\label{subsec_diversificacao_e_fronteira_eficiente}

O propósito central da diversificação de portfólios consiste em estruturar uma combinação de ativos que apresentem correlação nula ou moderada entre si, com a finalidade de maximizar os retornos ajustados ao risco do investimento \cite{bakry2021bitcoin}. A diversificação atua não apenas como um mecanismo de mitigação da exposição ao risco, mas também como um impulsionador para a obtenção de rentabilidade superior \cite{jeleskovic2024cryptocurrency}. Essa estratégia demonstra eficácia empírica mesmo quando incorpora classes de ativos caracterizadas por elevada volatilidade na composição da carteira \cite{bakry2021bitcoin}.

Sob a ótica do modelo de média-variância (MV), a formulação matemática configura-se como um problema de otimização multiobjetivo, cujo resultado é um conjunto de soluções ótimas denominado fronteira eficiente de investimentos \cite{paiva2019decision}. Essa fronteira delimita a carteira de ativos capaz de alcançar o patamar máximo de retorno esperado para um dado limite mínimo de risco estabelecido, indicando as estratégias de alocação mais eficientes para qualquer nível de retorno desejado pelo investidor \cite{paiva2019decision}.

A otimização de carteiras é a técnica voltada à extração do benefício máximo derivado da alocação de recursos financeiros, respeitando estritamente o perfil de risco-retorno do investidor \cite{jeleskovic2024cryptocurrency}. O procedimento requer o ajuste percentual do peso de cada ativo visando, em termos estatísticos, a minimização da variância ou desvio padrão e a maximização da média dos retornos \cite{jeleskovic2024cryptocurrency}. Profissionais de investimento exploram o potencial de diferentes ativos por meio de simulações e da aplicação de cenários de restrição sobre os pesos de alocação \cite{bakry2021bitcoin}.

Do ponto de vista empírico, a determinação da ponderação ideal frequentemente utiliza a maximização do Índice de Sharpe como objetivo principal, método que origina o que a literatura técnica denomina como carteira de tangência \cite{jeleskovic2024cryptocurrency}. As modelagens aplicam restrições orçamentárias unitárias, permitindo o posicionamento do investidor em ativos arbitrários dentro dos limites impostos pelo cenário \cite{bakry2021bitcoin}. Contudo, a aplicação dessa técnica em múltiplos ativos de alto risco pode gerar carteiras ótimas compostas por uma grande quantidade de frações diminutas (pesos de pequena escala), impondo barreiras práticas para a implementação da estratégia, especialmente no caso de investidores individuais \cite{chaweewanchon2022markowitz}.

\subsection{Aplicação da média-variância em criptoativos}
\label{subsec_aplicacao_da_media_variancia_em_criptoativos}

A integração crescente das criptomoedas na alocação de capital impõe a verificação sobre os reais benefícios de otimização que esses ativos conferem às carteiras \cite{jeleskovic2024cryptocurrency}. A característica de descorrelação do Bitcoin em relação aos ativos financeiros convencionais sinaliza uma vantagem potencial de diversificação quando adicionado a portfólios previamente diversificados \cite{bakry2021bitcoin}. Estudos reportam que a inserção de Bitcoin, mesmo em proporções marginais, produz resultados na relação média-variância superiores aos de carteiras compostas estritamente por ativos tradicionais \cite{bakry2021bitcoin}. A criptomoeda demonstra capacidade de gerar retornos ajustados ao risco em carteiras mistas de ações e títulos corporativos, apresentando propriedades de proteção \textit{(safe haven)} superiores às do ouro e de outras commodities \cite{bakry2021bitcoin}.

Análises evidenciam que portfólios híbridos, compostos por ativos convencionais e criptomoedas, exibem um Índice de Sharpe superior em avaliações ex-post e comportamentos mais estáveis em projeções ex-ante \cite{jeleskovic2024cryptocurrency}. Além das melhorias de desempenho estatístico, a alocação em criptoativos funciona como uma camada de proteção adicional para o investidor frente aos riscos atrelados a políticas monetárias que afetam o mercado de moedas fiduciárias tradicionais \cite{jeleskovic2024cryptocurrency}.
